% LaTeX template for Finance summaries
% Stu Alden
%----------------------------------------------------------------------------------------
% Simple Annuities


\documentclass[14pt]{extarticle}

\usepackage{conceptsheet}

%---------------------------------------
\begin{document}

\finmathtitle{4. Simple Annuities}

\textbf{Annuities} are simply a collection of multiple payments, generally (but not always)
level amounts that are generally (but not always) equally spaced.  When amounts are level and equally
spaced, and the rate period coincides with the payment period, we call them ``simple'' annuities.

The mathematics needed to value them is exactly the same as what we've already seen for compound interest,
but when there are many payments, it is usually easier to use a calculator or computer.  There are some
computational shortcuts, however, and that is the subject of this concept sheet.  We'll use the following
notation:

\begin{center}
  \renewcommand{\arraystretch}{1.2}
  \begin{tabular}{rl}
    \textbf{R} & Amount of each payment, paid at the \underline{end} of each period \\
    \textbf{n} & Number of payments, so here also the number of periods             \\
    \textbf{i} & Interest rate per period (as a decimal, \%/100)                    \\
    \textbf{A} & Present Value of the annuity (at time 0)                           \\
    \textbf{S} & Accumulated (Future) Value of the annuity (at time n)              \\
  \end{tabular}
\end{center}

Then the timeline looks like this:

\begin{center}
  \begin{tikzpicture}
    %draw horizontal line with squiggle in middle
    \draw (0,0) -- (5,0);
    \draw[decorate,decoration={snake,pre length=5mm, post length=5mm}] (5,0) -- (7,0);
    \draw (7,0) -- (10,0);

    %draw vertical tick-marks
    \foreach \x in {0,2,4,8,10}
    \draw (\x cm,3pt) -- (\x cm,-3pt);

    %Label "nodes"
    \draw (0,0) node[below=3pt] {$ 0 $} node[above=3pt] {$  $};
    \draw (0,0) node[below=18pt] {$ A $} node[above=3pt] {$  $};

    \draw (2,0) node[below=3pt] {$ 1 $} node[above=3pt] {$ R $};
    \draw (4,0) node[below=3pt] {$ 2 $} node[above=3pt] {$ R $};
    \draw (6,0) node[below=3pt] {$  $} node[above=3pt] {$  $};
    \draw (8,0) node[below=3pt] {$ n - 1 $} node[above=3pt] {$ R $};
    \draw (10,0) node[below=3pt] {$ n $} node[above=3pt] {$ R $};
    \draw (10,0) node[below=18pt] {$ S $} node[above=3pt] {$ R $};
  \end{tikzpicture}
\end{center}

Then we have

\begin{align*}
  S & = R \cdot \frac{(1 + i)^n-1}{i} \\
\end{align*}

\vspace{-.2in}

S denotes the ending, future, or \textbf{accumulated value} of the stream of annuity payments.
\pagebreak

To value the payments at the beginning instead (denoted \textbf{present value}), we just discount S by
n periods, as follows:

\begin{align*}
  A & = \frac{S}{(1 + i)^n}              \\\\
  & = R \cdot \frac{1-(1 + i)^{-n}}{i} \\\\
  & = R \cdot \frac{1 - v^n}{i}
\end{align*}

where, as before,

\[
  v = \frac{1}{1+i} = (1+i)^{-1}
\]

\vspace{.5in}

\textbf{Note:}

When actuaries want to specify the interest rate and the term of the annuity in the notation, they use the symbol
{\Large
  \[
    \ax{\angln i}
  \]
}
for present value and
{\Large
  \[
    \sx{ \angln i}
  \]
}
\vspace{.2in}

for accumulated or future value, where $ n $ is the number of periods and $ i $ is the rate per period.

\end{document}
